\section{Portability Model}
\label{sec:navhal_portability_model}

NavHAL achieves portability through structured separation of hardware concerns and compile-time specialization. Instead of relying on runtime abstraction layers, the system follows a clear adaptation flow where only the required components are implemented or replaced based on the target platform.

\subsection*{Porting Flow Overview}

Portability in NavHAL follows a layered adaptation path:

\begin{center}
\texttt{Application → Interface → Core → Vendor → Board → Hardware}
\end{center}

\begin{figure}[H]
    \centering
    \begin{tikzpicture}[
        node distance=0.5cm,
        layer/.style={draw, rectangle, minimum width=6.5cm, minimum height=0.7cm, align=center, font=\small, rounded corners=2pt},
        hwlayer/.style={draw, rectangle, fill=gray!10, minimum width=6.5cm, minimum height=0.7cm, align=center, font=\small\bfseries, rounded corners=2pt},
        arrow/.style={-Latex, thick}
    ]

    \node[layer, fill=blue!5] (app) {Application Logic (Hardware Independent)};
    \node[layer, below=of app] (api) {NavHAL Unified Interface (APIs)};
    
    \node[layer, below=1cm of api] (core) {\textbf{Core Layer} (Arch-specific Behavior)};
    \node[layer, below=of core] (vendor) {\textbf{Vendor Layer} (MCU Definitions)};
    \node[layer, below=of vendor] (board) {\textbf{Board Layer} (Pin Configuration)};
    \node[hwlayer, below=of board] (hw) {Physical Hardware (Silicon)};

    \draw[arrow] (app) -- (api);
    \draw[arrow] (api) -- (core);
    \draw[arrow] (core) -- (vendor);
    \draw[arrow] (vendor) -- (board);
    \draw[arrow] (board) -- (hw);

    % Cases
    \node[anchor=west, font=\tiny\bfseries\color{primary}] at ($(board.west)+(-4.5,0)$) {Case 1: Board Update};
    \node[anchor=west, font=\tiny\bfseries\color{primary}] at ($(vendor.west)+(-4.5,0)$) {Case 2: MCU Update};
    \node[anchor=west, font=\tiny\bfseries\color{primary}] at ($(core.west)+(-4.5,0)$) {Case 3: Arch Update};

    % Adaptation Zone
    \draw[decorate, decoration={brace, amplitude=10pt, raise=15pt}, secondary]
        (core.north east) -- (board.south east)
        node[midway, right=25pt, font=\small\bfseries, text width=2.5cm] {Adaptation Zone};

    \end{tikzpicture}
    \caption{Portability model of NavHAL illustrating the vertical adaptation stack. Only the lower layers are modified when porting between boards, MCUs, or architectures.}
    \label{fig:navhal_portability_model}
\end{figure}

The application and interface layers remain unchanged. Porting is performed by adapting the lower layers depending on what changes in the hardware platform.

\subsection*{Case 1: Porting to a New Board (Same MCU)}

When the microcontroller remains the same but the physical board changes, only the \textbf{board layer} needs to be updated.

\textbf{Flow:}
\begin{itemize}
    \item Define board-specific pin mappings (e.g., mapping logical pins to actual GPIOs)
    \item Configure default peripherals (LEDs, communication ports, etc.)
    \item Update linker script if memory layout differs
\end{itemize}

\textbf{Impact:}
\begin{itemize}
    \item No changes required in core or vendor layers
    \item No changes required in application code
\end{itemize}

This is the simplest portability case and enables reuse of the entire software stack across different hardware layouts.

\subsection*{Case 2: Porting to a New Microcontroller (Same Architecture)}

When moving to a different microcontroller within the same architecture (e.g., STM32F401 → STM32F411), the \textbf{vendor layer} is adapted.

\textbf{Flow:}
\begin{itemize}
    \item Define register base addresses and peripheral mappings
    \item Update memory layout and peripheral availability
    \item Provide vendor-specific definitions required by the core layer
\end{itemize}

The core layer continues to use these definitions to implement peripheral behavior.

\textbf{Impact:}
\begin{itemize}
    \item Core logic remains unchanged
    \item Board layer may require minor updates
    \item Application code remains unchanged
\end{itemize}

This separation allows reuse of peripheral logic while adapting only low-level hardware details.

\subsection*{Case 3: Porting to a New Architecture}

When targeting a completely different processor architecture (e.g., moving from Cortex-M4 to another architecture), a new \textbf{core layer} must be implemented.

\textbf{Flow:}
\begin{itemize}
    \item Implement architecture-specific startup code and interrupt handling
    \item Provide core implementations for all common HAL APIs
    \item Integrate vendor-specific register definitions for the new platform
\end{itemize}

\textbf{Impact:}
\begin{itemize}
    \item Vendor and board layers are built on top of the new core
    \item Interface remains unchanged
    \item Application code remains unchanged
\end{itemize}

Although this is the most involved step, it is performed once per architecture and reused across multiple platforms.

\subsection*{Compile-Time Binding}

All portability decisions are resolved at compile time through build configuration parameters such as:

\begin{itemize}
    \item Target processor (core)
    \item Vendor family
    \item Board configuration
\end{itemize}

The build system includes only the selected implementations, ensuring:

\begin{itemize}
    \item No runtime branching or hardware detection
    \item Reduced binary size
    \item Direct mapping to hardware
\end{itemize}

\subsection*{Key Observation}

In all cases, the upper layers of the system remain unchanged. Portability is achieved by replacing only the minimal required layer:

\begin{center}
\begin{tabular}{|c|c|}
\hline
\textbf{Change Type} & \textbf{Layer Modified} \\
\hline
New Board & Board Layer \\
\hline
New MCU (same arch) & Vendor Layer \\
\hline
New Architecture & Core Layer \\
\hline
\end{tabular}
\end{center}